% !TEX root = ../main.tex
\section{Introduction}
\label{sec:intro}

Large language model (LLM) agents increasingly rely on external \emph{skills}---structured tool descriptions that extend model capabilities beyond parametric knowledge---to perform domain-specific tasks.
Recent work demonstrates that curated skill libraries substantially improve agent performance: \citet{skillsbench2026} report an average gain of 16.2 percentage points across domains, with the medical domain exhibiting the largest improvement at 51.9 percentage points.
This finding underscores the critical role of skill quality in high-stakes settings such as clinical decision support, biomedical literature analysis, and drug safety assessment.

Despite this promise, three gaps hinder progress toward reliable medical agent skills.

\paragraph{Gap 1: Task-level evaluation obscures skill-level failures.}
Existing medical agent benchmarks evaluate end-to-end task success \citep{healthbench2025, medagentbench2025, labbench2024} but do not isolate whether the agent correctly \emph{executed} the skill specification it was given.
An agent may answer a question correctly by ignoring a poorly written skill description entirely, or it may fail a task despite faithfully following a description that lacks critical information.
Task-level metrics conflate these fundamentally different failure modes.

\paragraph{Gap 2: Skills are treated as static artifacts.}
Current skill development follows a ``write once, deploy forever'' paradigm \citep{skillnet2026, easytool2024}.
Once a \skillmd file is authored, it is rarely revisited---even when downstream execution data reveals systematic failure patterns.
This stands in contrast to conventional software engineering, where documentation is continuously refined based on usage feedback.

\paragraph{Gap 3: Description optimization at scale remains unexplored.}
Prompt evolution methods \citep{gepa2025, play2prompt2025} have shown that iterative refinement of natural-language specifications can improve LLM performance without weight updates.
However, these techniques have not been applied to large-scale skill libraries, particularly in the medical domain where description precision directly affects patient safety.

\medskille addresses these gaps through an \textbf{evaluate--diagnose--evolve} closed loop that treats skill descriptions as first-class, evolvable agent interfaces rather than frozen documentation.\footnote{Code, data, and evolved skill descriptions are available at \anonrepo.}

\begin{figure}[t]
\figbox{5.2cm}{Figure 1 (Teaser): the \rbu gap. Left panel---agents read \skillmd at 96.2\% invocation. Right panel---only 74.3\% complete the specified task. The 22-pp wedge between the two represents specification grounding failure that \gepamed iteratively closes.}
\caption{The \rbu gap: agents read skill specifications (96.2\%) yet only complete skill-grounded tasks at 74.3\%, exposing a 22-percentage-point specification grounding failure that \gepamed targets via reflective description evolution.}
\label{fig:teaser}
\end{figure}

\paragraph{Contribution 1: \medskillbench.}
We introduce the first large-scale skill-level execution benchmark for medical AI agents.
\medskillbench aggregates 1,462 skills from three open-source medical skill repositories and evaluates each in an isolated sandbox environment.
A mandatory \emph{read-first} protocol ensures the agent must parse the \skillmd specification before attempting execution.
An LLM-as-judge scores each execution trajectory along eight binary dimensions---\emph{skill invocation}, \emph{correct usage}, \emph{task completion}, \emph{tool execution success}, \emph{output quality}, \emph{reasoning quality}, \emph{efficiency}, and \emph{trajectory quality}---providing fine-grained diagnostic signal rather than a single pass/fail label.

\paragraph{Contribution 2: \rbu gap quantification.}
Using \medskillbench, we identify and quantify a systematic \emph{read-but-not-use} (\rbu) gap: agents invoke skills at a rate of 96.2\% but complete the specified tasks at only 74.3\%, yielding a 22-percentage-point gap.
This gap represents a \emph{specification grounding failure}---the model recognizes what a skill describes but cannot reliably translate that description into correct execution.
Stratifying by \texttt{tool\_type} reveals that the gap is strongly modulated by skill interface design: guide-style skills (no executable tool) exhibit an \rbu gap of 0.269, while tool-based skills show substantially smaller gaps (CLI: 0.106; Python: 0.099; mixed: 0.069; R: 0.111).
This finding shifts the diagnosis from ``models cannot use tools'' to ``models struggle most when skill descriptions lack executable grounding.''

\paragraph{Contribution 3: \gepamed.}
We propose \gepamed, a reflective prompt evolution method adapted from GEPA \citep{gepa2025} for medical \skillmd optimization.
\gepamed uses execution trajectories and judge feedback from \medskillbench as evolutionary signal to iteratively rewrite skill descriptions---improving groundability without any model weight changes.
\todo{GEPA-Med experimental results pending; methodology described in \S\ref{subsec:gepa}.}

Our work is organized around four research questions:
\begin{enumerate}[nosep,leftmargin=*,label=\textbf{RQ\arabic*:}]
    \item Do medical agent skills exhibit a systematic read-but-not-use gap between specification comprehension and execution?
    \item Is the gap primarily attributable to environment dependency or to description-level grounding failure?
    \item Can description-only evolution (no weight changes) improve skill groundability?
    \item Do evolved skill descriptions transfer across different LLM backends?
\end{enumerate}

\paragraph{Positioning.}
Following the taxonomy of \citet{skillsbench2026}, \medskille operates primarily at the \emph{Skill Deployment} tier (T2), evaluating and improving skills post-authoring.
It secondarily informs \emph{Skill Development} (T1) by identifying which description patterns yield higher groundability.
Unlike reinforcement learning approaches to post-deployment optimization \citep{tracefree2026}, \medskille is entirely training-free: it modifies only the natural-language skill description, making it applicable to closed-weight models and resource-constrained settings alike.

The remainder of this paper is structured as follows.
\S\ref{sec:related_work} reviews related work.
\S\ref{sec:method} details our framework, including the \medskillbench evaluation protocol and the \gepamed evolution methodology.
\S\ref{sec:experiments} reports experimental results and discusses implications.
\S\ref{sec:conclusion} concludes.
